% Paper 2 — lot release attributes and reactogenicity of VA-MENGOC-BC (English).
% Target journal: Vaccine (Elsevier; research article, abstract <= 250 words).
\documentclass[11pt]{article}
\newcommand{\DocLang}{en}
\input{papers/shared/preamble.tex}
\input{papers/shared/authors.tex}
\input{papers/shared/pending.tex}
\input{papers/p2-lot-quality/shared/flow.tex}
\externaldocument{supplement}
\usepackage{lineno}
\onehalfspacing

\hypersetup{pdftitle={Lot release attributes and reactogenicity of the VA-MENGOC-BC outer membrane vesicle vaccine},
  pdfauthor={Richard Alejandro Matos Arderi; Abel Ponce Gonzalez}}

\begin{document}

\begin{titlepage}
\SyntheticNotice
\begin{center}
{\LARGE\bfseries Lot release attributes and reactogenicity of an outer membrane vesicle vaccine:
linkage of quality-control data from \PTNLots{} VA-MENGOC-BC lots with national adverse event
reports\par}
\bigskip
\AuthorBlock
\end{center}
\bigskip
\noindent\textbf{Running title:} Lot quality and reactogenicity of VA-MENGOC-BC\par
\noindent\textbf{Keywords:} vaccine lot release; quality control; statistical process control;
outer membrane vesicle; endotoxin; adverse events following immunization; case-only design;
pharmacovigilance\par
\end{titlepage}

\linenumbers

\section*{Abstract}
\noindent\textbf{Background:} Whether the measured quality of vaccine lots influences reported
reactions is debated and unstudied for outer membrane vesicle (OMV) vaccines, whose reactogenicity
depends on endotoxin and adjuvant.\par
\noindent\textbf{Methods:} Release data of \PTNLots{} VA-MENGOC-BC lots (\PTLotFirstYear--\PTLotLastYear;
\PTNAttributes{} attributes) were analysed with process control, capability ($P_{pk}$), robust
multivariate monitoring, change-point detection and national-versus-export equivalence tests.
Cuban reports of adverse events following immunization (\VStudyFirstYear--\VStudyLastYear) were
linked to the administered lots. In a case-only design, generalised estimating equations clustered
by lot estimated whether reports of fever or severe local reactions varied with lot attributes, with
false-discovery-rate control, a negative-control exposure and sensitivity analyses.\par
\noindent\textbf{Results:} \PTConformitySentence; $P_{pk}$ was below 1.0 for \PTNAttrPpkBelowOne{}
attributes and national and export lots were equivalent for \PTNEquivalent{} of \PTNAttributes.
Of \PTNTargetReports{} reports, \PTNLinked{} (\PTLinkRate) were linked to \PTNLinkedLots{} lots.
\PTNSignificantFDR{} of \PTNPrimaryTests{} prespecified associations survived correction
(\PTSignificantList); the odds ratio per standard deviation of endotoxin for fever of at least
39\,\textdegree{}C was \PTEvFeverThreeNineEndotoxinOR{} (\CI{} \PTEvFeverThreeNineEndotoxinCI).
\PTNegControlSentence, and lot attributes \PTIncrementalWord{} out-of-lot prediction.\par
\noindent\textbf{Conclusions:} Linking routine release testing with pharmacovigilance is feasible
and tests lot-related hypotheses about reactogenicity within specification, supporting quality risk
management and post-marketing surveillance of OMV vaccines.

\clearpage
\section{Introduction}\label{sec:intro}

Every lot of a vaccine is released only after its quality attributes have been tested against the
specifications of its licence \cite{WHO2013LotRelease}. Within those limits, attributes such as
antigen content, adsorption to the adjuvant or residual endotoxin still vary from lot to lot. Whether
this variability within specification influences the reactions that vaccinees experience is a
question that pharmaceutical quality systems \cite{ICHQ9R1,ICHQ10} and pharmacovigilance rarely
answer together. Interest in the question has grown since analyses of national reporting systems
suggested heterogeneity in suspected adverse events across lots of COVID-19 vaccines
\cite{Schmeling2023,Manniche2024,Furst2024,Schmeling2025}. Those analyses also showed the pitfalls
of the approach: the number of doses actually administered from each lot, differences in the
populations vaccinated and changes in reporting over time can create apparent differences between
lots \cite{Hviid2023,Scott2023}. Methodological work has proposed scan statistics over the
manufacturing genealogy to locate such signals \cite{Mahaux2019}, but no study has related measured
release attributes to reported reactions.

Outer membrane vesicle (OMV) vaccines against serogroup~B meningococcal disease are a natural
setting for this question. Their reactogenicity---mainly fever and local reactions in young
children---is well described \cite{Nokleby2007,Holst2009}, and it has plausible lot-level
determinants: the membrane-bound endotoxin of the vesicles drives the febrile response
\cite{Sheerin2019,Zariri2016}, and adsorption onto aluminium hydroxide reduces pyrogenicity
\cite{Rosenqvist1998}. VA-MENGOC-BC, the Cuban OMV vaccine against serogroups~B and~C, has been
produced continuously for more than three decades and is administered to every Cuban infant
\cite{Sotolongo2007,SierraGonzalez2019}. Its manufacturer keeps the release record of every lot,
and the national surveillance system records the lot administered in each report of an adverse
event following immunization (AEFI) \cite{Galindo2012}.

We linked these two sources with two aims: (i) to characterise the consistency of VA-MENGOC-BC
production with statistical process control, process capability and multivariate monitoring; and
(ii) to estimate, in a case-only design, whether lot attributes are associated with the profile of
reactions reported after vaccination, with prespecified hypotheses, a negative-control exposure
and an explicit assessment of statistical power.

\section{Methods}\label{sec:methods}

\subsection{Design and data sources}\label{sec:design}
The study combined a retrospective analysis of lot release data with a case-only analysis of
linked AEFI reports, and is reported following RECORD-PE \cite{Langan2018}. The release record
contained \PTNLots{} lots produced in \PTLotFirstYear--\PTLotLastYear, \PTNNationalLots{} for the
national programme and \PTNExportLots{} for export, with \PTNAttributes{} quality attributes and
their specifications. The AEFI reports came from the national surveillance database described in a
companion analysis: \VNFiles{} annual files of \VStudyFirstYear--\VStudyLastYear, harmonised,
pseudonymised and deduplicated with the same open-source pipeline (version \VCodeVersion)
\cite{Kahn2016,ENISA2019}. \VFormSentence. Lot data are confidential industrial information; they were analysed
on an institutional workstation and are released only as positions within the specification
window (0 = lower limit, 1 = upper limit), never as raw values.

\subsection{Attributes and scales}\label{sec:scales}
Attributes were analysed on their natural scale or after the transformation that makes their
variation approximately symmetric: base-2 logarithm for the bactericidal titre, base-10 logarithm
for the anti-OMV IgG concentration and $\log_{10}(1+x)$ for endotoxin. For the association
analyses, each attribute was standardised to the mean and standard deviation (SD) of the linked
lots, so that odds ratios refer to a one-SD difference between lots.

\subsection{Process performance}\label{sec:m-spc}
For each attribute in production order we built individual and moving-range charts with Nelson's
rules \cite{Nelson1984,Montgomery2019} and an exponentially weighted moving average chart
($\lambda=\PTEWMALambda$) \cite{Lucas1990}. Process performance was summarised with $P_{pk}$
with \PTCapBootstrap{} percentile bootstrap intervals \cite{Efron1987} and with the
percentile-based $P_{pk}$ of ISO~22514-2, which does not assume normality \cite{ISO22514-2}, for
all lots and for lots produced before (\PTPeriodOne) and during (\PTPeriodTwo) the
pharmacovigilance window; the split was fixed by the study design, not by the data. Abrupt shifts
were located with an exact PELT algorithm with a BIC-type penalty \cite{Killick2012,Truong2020}.
Multivariate atypicality was monitored with a robust principal component model (minimum covariance
determinant \cite{Rousseeuw1999}, components explaining at least 85\% of the variance), with
Hotelling's $T^2$ control limits from the beta distribution \cite{Tracy1992} and the squared
prediction error with Jackson--Mudholkar limits \cite{Jackson1979,MacGregor1995}, both at
$\alpha=\PTMSPCAlpha$. National and export lots were compared with two one-sided tests
\cite{Schuirmann1987} (Welch standard errors; equivalence margin of \PTEqMarginPct{} of the
specification window), the Hodges--Lehmann shift \cite{Hodges1963} and a multivariate energy test
\cite{Szekely2013}.

\subsection{Linkage of reports to lots}\label{sec:m-link}
Lot numbers recorded in the reports were normalised and matched to the release record in three
steps: exact match, match on the core of the lot code (ignoring separators and suffixes), and
fuzzy match (similarity of at least \PTFuzzyThreshold{} among national lots produced before
vaccination). Ambiguous matches were resolved in favour of national lots and of the most recent
lot produced before the year of vaccination. Links implying vaccination more than
\PTMaxYearsAfterProduction{} years after production were considered implausible and excluded from
the primary analysis (\figref{fig:flow}). Characteristics of linked and unlinked reports were
compared to assess linkage bias (Supplementary \tabref{tab:S-linkbias}).

\begin{figure}[tbp]
\centering
\LinkageFlow
\caption{Linkage of AEFI reports involving VA-MENGOC-BC to the release record of the administered
lots. Counts come from the disclosure-controlled release.}
\label{fig:flow}
\end{figure}

\subsection{Case-only association analysis}\label{sec:m-assoc}
Because the number of doses administered from each lot is not recorded, absolute risks per lot
cannot be estimated. We therefore used a case-only design \cite{Piegorsch1994}: among reports
linked to a lot, we estimated whether the probability that a report included a given reaction
varied with the attributes of the lot. The design is valid for the relative profile of reactions
under the assumption that the completeness of reporting of one reaction relative to another does
not vary with the attributes of the lot.

The prespecified primary family comprised three outcomes (fever of at least 39 and
40\,\textdegree{}C and severe local reaction) and four exposures with a plausible mechanism
(endotoxin, aluminium hydroxide, protein concentration and protein adsorption), that is
\PTNPrimaryTests{} associations. Fill volume, which should not influence the type of reaction, was
a negative-control exposure \cite{Lipsitch2010,Arnold2016}. Each association was estimated with a
logistic generalised estimating equation with exchangeable working correlation within lot and
robust standard errors \cite{Liang1986,Zeger1986}, adjusted for age (log months), sex, dose number,
co-administration and calendar year; unadjusted estimates are reported alongside. The Benjamini--Hochberg procedure controlled the false
discovery rate at \PTFDRAlpha{} within the primary family \cite{Benjamini1995}; the negative control
was not part of the family. The other attributes, the multivariate atypicality ($T^2$) and the
secondary outcomes were exploratory.

Sensitivity analyses restricted the data to exact links, excluded co-administered vaccines and
retained temporally implausible links; refitted the exposures in a Bayesian random-intercept
logistic model; and modelled the proportion of reports with each reaction at lot level with a
quasi-binomial model weighted by the number of reports per lot. To interpret null findings, the
minimum detectable odds ratio per SD with 80\% power was computed from the number of reports, the
prevalence of each reaction and the intra-lot correlation (design effect) \cite{Hsieh1998}.

\subsection{Incremental predictive value}\label{sec:m-ml}
We asked whether the attributes of a lot help predict the reactions reported with it beyond the
characteristics of the vaccinee. Gradient-boosted trees \cite{Ke2017} were trained with and without
the lot attributes and evaluated by the area under the ROC curve in \PTCVFolds-fold cross-validation
grouped by lot, so that every prediction concerns lots unseen during training. The difference in
area was compared with its distribution under \PTNPermutations{} permutations of complete attribute
profiles between lots, which preserves the correlation between attributes and the clustering of
reports.

\subsection{Disclosure control and ethics}\label{sec:m-sdc}
Only aggregated results left the controlled environment: non-zero counts below \VMinCell{} were
suppressed together with the statistics derived from them, lot attributes were released only as
positions within the specification window, and an automated verifier checked every released file
\cite{Hundepool2012}. The study used pseudonymised surveillance records and manufacturing records
under the institutional data-governance procedure of the Finlay Vaccine Institute and Law 149/2022
\cite{Ley149}.

\section{Results}\label{sec:results}

\subsection{Consistency of production}\label{sec:r-spc}
\PTConformitySentence{} (lowest percentage of conforming lots for any attribute,
\PTPctConforming). $P_{pk}$ was at least 1.33 for \PTNAttrPpkAbove{} attributes and
below 1.0 for \PTNAttrPpkBelowOne{} (\PTPpkBelowOneList; \tabref{tab:capability},
\figref{fig:capability}); values below 1.0 indicate that the process is centred or spread such
that lots close to a limit are expected. The PELT algorithm located
\PTNChangepoints{} sustained shifts (\PTChangepointList). The robust principal component model
retained \PTNComponents{} components explaining \PTExplainedVar{} of the variance; \PTTTwoFlagged{}
lots exceeded the $T^2$ limit and \PTSPEFlagged{} the squared-prediction-error limit
(\figref{fig:t2}). National and export lots were equivalent within the prespecified margin for
\PTNEquivalent{} of \PTNAttributes{} attributes (Supplementary \tabref{tab:S-equivalence}), and their
multivariate distributions \PTEnergyWord{} (energy test $p$\PTEnergyP).

\begin{table}[tbp]
\centering\small
\caption{Process performance of VA-MENGOC-BC lots: $P_{pk}$ with 95\% bootstrap interval for all
lots and by production period, and production year of the change points located by PELT.}
\label{tab:capability}
\PTTableCapability
\end{table}

\begin{figure}[tbp]
\centering
\includegraphics[width=\textwidth]{p2_qc_annual}
\caption{Release attributes of VA-MENGOC-BC lots by production year, shown as position within the
specification window (0, lower limit; 1, upper limit). Lines show the annual median and band the
interquartile range; red lines mark specification limits.}
\label{fig:qcannual}
\end{figure}

\begin{figure}[tbp]
\centering
\begin{minipage}[t]{0.49\textwidth}\centering
\includegraphics[width=\linewidth]{p2_capability}
\caption{$P_{pk}$ (95\% bootstrap interval) by production period. Vertical lines mark 1.00 and
1.33.}
\label{fig:capability}
\end{minipage}\hfill
\begin{minipage}[t]{0.49\textwidth}\centering
\includegraphics[width=\linewidth]{p2_t2}
\caption{Multivariate atypicality ($T^2$) of lots by production year, with the control limit.}
\label{fig:t2}
\end{minipage}
\end{figure}

\subsection{Linkage}\label{sec:r-link}
Of \PTNTargetReports{} reports involving VA-MENGOC-BC, \PTNLinked{} (\PTLinkRate) were linked to
\PTNLinkedLots{} lots: \PTLinkExact{} by exact match, \PTLinkCore{} by the core of the code and
\PTLinkFuzzy{} by fuzzy match (\figref{fig:flow}; Supplementary \tabref{tab:S-linkage}). The primary
analysis included \PTNAnalysed{} reports in \PTNAnalysedLots{} lots (median \PTReportsPerLot{}
reports per lot).

\subsection{Lot attributes and reported reactions}\label{sec:r-assoc}
\PTNSignificantFDR{} of the \PTNPrimaryTests{} prespecified associations remained significant after
false-discovery-rate correction (\PTSignificantList; \tabref{tab:gee}, \figref{fig:forest}). The
proportion of reports with fever of at least 39\,\textdegree{}C was
\PTEvFeverThreeNineEndotoxinDir{} the endotoxin content of the lot (odds ratio per SD
\PTEvFeverThreeNineEndotoxinOR, \CI{} \PTEvFeverThreeNineEndotoxinCI; $q$\PTEvFeverThreeNineEndotoxinQ)
and \PTEvFeverThreeNineAlohThreeConcDir{} its aluminium hydroxide content (\PTEvFeverThreeNineAlohThreeConcOR,
\CI{} \PTEvFeverThreeNineAlohThreeConcCI). Fever of at least 40\,\textdegree{}C was
\PTEvFeverFourZeroEndotoxinDir{} endotoxin (\PTEvFeverFourZeroEndotoxinOR, \CI{}
\PTEvFeverFourZeroEndotoxinCI), and severe local reactions were \PTEvSevereLocalReactionAlohThreeConcDir{}
aluminium hydroxide (\PTEvSevereLocalReactionAlohThreeConcOR, \CI{}
\PTEvSevereLocalReactionAlohThreeConcCI). \PTNegControlSentence{} ($p<0.05$ in \PTNNegControlNominal{} of three outcomes). \PTSensSentence{} (Supplementary
\tabref{tab:S-sensitivity}).

\begin{table}[tbp]
\centering\small
\caption{Association between lot attributes and the reactions reported with VA-MENGOC-BC
(case-only design). Odds ratios per SD of the attribute between lots, from logistic generalised
estimating equations clustered by lot, crude and adjusted for age, sex, dose, co-administration and
year;
$q$, Benjamini--Hochberg adjusted $p$ within the prespecified family (the negative control is
outside the family).}
\label{tab:gee}
\PTTableGEE
\end{table}

\begin{figure}[tbp]
\centering
\includegraphics[width=\textwidth]{p2_forest}
\caption{Adjusted odds ratios per SD of each lot attribute (95\% CI) for each reaction,
single-exposure models. The shaded band spans the minimum detectable effect (from 1/MDE to MDE);
grey markers show the negative-control exposure.}
\label{fig:forest}
\end{figure}

\subsection{Power and incremental predictive value}\label{sec:r-power}
With the observed number of lots and intra-lot correlation, the minimum detectable odds ratio per
SD was \PTEvFeverThreeNineMDE{} for fever of at least 39\,\textdegree{}C, \PTEvFeverFourZeroMDE{}
for fever of at least 40\,\textdegree{}C and \PTEvSevereLocalReactionMDE{} for severe local
reactions (\tabref{tab:power}). Adding all lot attributes to the characteristics of the vaccinee
\PTIncrementalWord{} the out-of-lot prediction of the reactions: the change in the area under the
ROC curve was \PTEvFeverThreeNineDeltaAUC{} for fever of at least 39\,\textdegree{}C (permutation
$p$\PTEvFeverThreeNinePermP; \figref{fig:incremental}).

\begin{table}[tbp]
\centering\small
\caption{Minimum detectable effect (MDE; odds ratio per SD with 80\% power, two-sided
$\alpha=0.05$) given the clustering of reports in lots, and incremental value of the lot attributes
for out-of-lot prediction ($\Delta$AUC, change in the area under the ROC curve; $p$ from
permutations of attribute profiles between lots).}
\label{tab:power}
\PTTablePower
\end{table}

\begin{figure}[tbp]
\centering
\includegraphics{p2_incremental}
\caption{Change in the cross-validated area under the ROC curve when lot attributes are added
(points), against its permutation distribution (95th percentile marked).}
\label{fig:incremental}
\end{figure}

\section{Discussion}\label{sec:discussion}

\subsection{Principal findings}\label{sec:d-principal}
The release record of \PTNLots{} lots of VA-MENGOC-BC showed a process
\PTWithinSpecPhrase, with \PTNChangepoints{} sustained shifts and capability below 1.0 for
\PTNAttrPpkBelowOne{} attributes. Linking \PTLinkRate{} of the AEFI reports to their lots, we found
that \PTNSignificantFDR{} of \PTNPrimaryTests{} prespecified associations between lot attributes
and the profile of reported reactions survived correction for multiplicity (\PTSignificantList),
while the negative-control exposure \PTNegControlWord{} with the reactions.

\subsection{Interpretation}\label{sec:d-interpretation}
An association between endotoxin and fever would be biologically coherent: OMV-associated
lipooligosaccharide activates innate immunity through Toll-like receptor~4, and the febrile
response to OMV vaccines has been attributed to it \cite{Sheerin2019,Zariri2016}, while adsorption
onto aluminium hydroxide modulates pyrogenicity \cite{Rosenqvist1998}. Any such association must be
read for what the design can show: a shift in the composition of the reports, among lots that were
all within specification, not an absolute risk per dose. The negative-control exposure, the
sensitivity analyses and the minimum detectable effects help separate true lot effects from the
artefacts that have affected lot-level analyses of spontaneous reports---differences in the
populations vaccinated with each lot, in calendar time and in reporting
\cite{Hviid2023,Scott2023}. For the same reason, an absence of association is informative only for
effects larger than the minimum detectable odds ratios reported here.

Statistical process control applied retrospectively to release data adds a different kind of
evidence: whether the process remained stable, whether it changed abruptly and whether lots for
national use and for export were equivalent. A capability index below 1.0 does not mean
non-conformity, but indicates that the margin between the process and its limits is narrow and is a
natural target for quality risk management \cite{ICHQ9R1,ICHQ10}. Combining these indicators with
the safety linkage turns the lot release record into a tool for post-marketing surveillance.

\subsection{Strengths and limitations}\label{sec:d-limitations}
The study links, for the first time for an OMV vaccine, the measured quality of each lot with the
reactions reported after its administration; it uses the complete release record and the national
reports of \VStudyFirstYear--\VStudyLastYear, prespecifies its hypotheses and controls multiplicity, and generates every number
from a verified, disclosure-controlled release.

It has important limitations. First, the doses administered per lot were unknown, so the design
estimates the relative profile of reactions, not their incidence; an attribute that raised the
frequency of all reactions equally would not be detected. Second, reporting completeness could
differ between lots in ways related to their attributes, for example through calendar time;
adjustment for year and the negative control mitigate but do not exclude this bias. Third, lot
numbers were missing or unmatched in some reports, and linked and unlinked reports may differ
(Supplementary \tabref{tab:S-linkbias}). Fourth, the lots were \PTWithinSpecPhrase, so the
results do not extrapolate beyond the specified range. Fifth, release tests
measure a sample of each lot at release and do not capture changes during storage and distribution.
Finally, reactions were recorded as indicators of the national form without clinical validation.

\subsection{Implications}\label{sec:d-implications}
Manufacturers already hold the release records and national systems already record the lot
administered; the pipeline used here shows that both can be linked routinely, with confidentiality
of industrial data preserved through disclosure control. Recording the doses distributed per lot
would allow incidence per lot to be estimated, and the same framework can be applied to other
vaccines with lot-level quality data.

\section{Conclusions}\label{sec:conclusions}
VA-MENGOC-BC lots were produced \PTWithinSpecPhrase. Linking lot release attributes
with national AEFI reports is feasible and provides a reproducible way to test whether variation
within specification is associated with reactogenicity; the associations identified here
(\PTSignificantList) define hypotheses for confirmation and for quality risk management.

\section*{Declarations}
\input{papers/p2-lot-quality/shared/declarations_en.tex}

\bibliographystyle{vancouver}
\bibliography{papers/shared/bib/pubmed,papers/shared/bib/manual}

\end{document}
